%% Follow comments to support use.

%%%%%%%%%%%%%%%%%%%%%%%%%%%%%%%%%%%%%%%%%%%%%%%%%%%%%%%%%
%% STEP 1: Choose options for MSc / BSc / seminar layout and your bibliographic style
%%%%%%%%%%%%%%%%%%%%%%%%%%%%%%%%%%%%%%%%%%%%%%%%%%%%%%%%%

%%  Language: 
%%      finnish, swedish, or english
%%  Pagination (use twoside by default)  
%%      oneside or twoside,
%%  Study programme / kind of report
%%      csm  = Master's thesis in Computer Science Master's Programme;
%%      tkt = Bachelor's thesis in Computer Science Bachelor's Programme;
%%      seminar = seminar report
%%  For Master's thesis choose your line or track:
%%      (30 cr thesis, 2020 onwards, Master's Programme in Computer Science = csm)
%%      software-track-2020 = Software study track
%%      algorithms-track-2020 = Algorithms study track
%%      networking-track-2020 = Networking study track

\documentclass[finnish,twoside,censored,tkt]{HYthesisML} 
%\documentclass[english,twoside,censored,tkt]{HYthesisML} 

% If wanted, open new chapters only at right page.
% By default, "openany".
%\PassOptionsToClass{openright,twoside,a4paper}{report}
\PassOptionsToClass{openany,twoside,a4paper}{report}

\usepackage{csquotes}
%%%%%%%%%%%%%%%%%%%%%%%%%%%%%%%%%%%%%%%%%%%%%%%%%%%%%%%%%
%% REFERENCES
%% Some notes on bibliography usage and options:
%% natbib -> you can use, e.g., \citep{} or \parencite{} for (Einstein, 1905); with APA \cite -> Einstein, 1905 without ()
%% maxcitenames=2 -> only 2 author names in text citations, if more -> et al. is used
%% maxbibnames=99 as no great need to suppress the biliography list in a thesis
%% for more information see biblatex package documentation, e.g., from https://ctan.org/pkg/biblatex 

%% Reference style: select one 
%% for APA = Harvard style = authoryear -> (Einstein, 1905) use:
\usepackage[style=authoryear,bibstyle=authoryear,backend=biber,natbib=true,maxnames=99,maxcitenames=2,uniquelist=minyear,giveninits=true,uniquename=mininit]{biblatex}
%% for numeric = Vancouver style -> [1] use:
%\usepackage[style=numeric,bibstyle=numeric,backend=biber,natbib=true,maxbibnames=99,giveninits=true,uniquename=init]{biblatex}


\addbibresource{bibliography.bib}
% in case you want the final delimiter between authors & -> (Einstein & Zweistein, 1905) 
% \renewcommand{\finalnamedelim}{ \& }
% List the authors in the Bibilipgraphy as Lastname F, Familyname G,
\DeclareNameAlias{sortname}{family-given}
% remove the punctuation between author names in Bibliography 
%\renewcommand{\revsdnamepunct}{ }


%% Block of definitions for fonts and packages for picture management.
%% In some systems, the figure packages may not be happy together.
%% Choose the ones you need.

%\usepackage[utf8]{inputenc} % For UTF8 support, in some systems. Use UTF8 when saving your file.

\usepackage{lmodern}         % Font package, again in some systems.
\usepackage{textcomp}        % Package for special symbols
\usepackage[pdftex]{color, graphicx} % For pdf output and jpg/png graphics
\usepackage{epsfig}
\usepackage{subfigure}
\usepackage[pdftex, plainpages=false]{hyperref} % For hyperlinks and pdf metadata
\usepackage{fancyhdr}        % For nicer page headers
\usepackage{tikz}            % For making vector graphics (hard to learn but powerful)
%\usepackage{wrapfig}        % For nice text-wrapping figures (use at own discretion)
\usepackage{amsmath, amssymb} % For better math

\singlespacing               %line spacing options; normally use single

\fussy
%\sloppy                      % sloppy and fussy commands can be used to avoid overlong text lines
% if you want to see which lines are too long or have too little stuff, comment out the following lines
% \overfullrule=1mm
% to see more info in the detailed log about under/overfull boxes...
% \showboxbreadth=50 
% \showboxdepth=50



%%%%%%%%%%%%%%%%%%%%%%%%%%%%%%%%%%%%%%%%%%%%%%%%%%%%%%%%%
%% STEP 2:
%%%%%%%%%%%%%%%%%%%%%%%%%%%%%%%%%%%%%%%%%%%%%%%%%%%%%%%%%
%% Set up personal information for the title page and the abstract form.
%% Replace parameters with your information.
\title{Title}

\author{Firstname Lastname}
\date{\today}

% Set supervisors, use the titles according to the thesis language
% in English Prof. or Dr., or in Finnish toht. or tri or FT, TkT, Ph.D. or in Swedish... 
% Include both advisors and supervisors in the same entry to get them appear on the abstract page.
% For MSc theses, the terms have somehow changed to supervisors and examiners, denoting BSc level advisors and supervisors. An unfortunate mismatch for now. 
% Advisors for BSc theses and supervisors for MSc theses are the first contacts, supervisors for BSc theses and examiners for MSc theses are less visible but are responsible on % grading the works. 

\supervisors{Prof.~D.U.~Mind, Dr.~O.~Why}

\keywords{algorithms, data structures}
\additionalinformation{\translate{\track}}

%% For seminar reports:
%%\additionalinformation{Name of the seminar}

%% Provide classification terms, to appear on the abstract page.
%% Replace the classification terms below with the ones that match your work.
%% ACM Digital library provides a taxonomy and a tool for classification
%% in computer science. Use 1-3 paths, and use right arrows between the
%% about three levels in the path; each path requires a new line.

\classification{\protect{\ \\
\  General and reference $\rightarrow$ Document types  $\rightarrow$ Surveys and overviews\  \\
\  Applied computing  $\rightarrow$ Document management and text processing  $\rightarrow$ Document management $\rightarrow$ Text editing
}}

%% If you want to quote someone special. You can comment this line out and there will be nothing on the document.
%\quoting{Bachelor's degrees make pretty good placemats if you get them laminated.}{Jeph Jacques}


%% OPTIONAL STEP: Set up properties and metadata for the pdf file that pdfLaTeX makes.
%% Your name, work title, and keywords are recommended.
\hypersetup{
    unicode=true,           % to show non-Latin characters in Acrobat’s bookmarks
    pdftoolbar=true,        % show Acrobat’s toolbar?
    pdfmenubar=true,        % show Acrobat’s menu?
    pdffitwindow=false,     % window fit to page when opened
    pdfstartview={FitH},    % fits the width of the page to the window
    pdftitle={},            % title
    pdfauthor={},           % author
    pdfsubject={},          % subject of the document
    pdfcreator={},          % creator of the document
    pdfproducer={pdfLaTeX}, % producer of the document
    pdfkeywords={something} {something else}, % list of keywords for
    pdfnewwindow=true,      % links in new window
    colorlinks=true,        % false: boxed links; true: colored links
    linkcolor=black,        % color of internal links
    citecolor=black,        % color of links to bibliography
    filecolor=magenta,      % color of file links
    urlcolor=cyan           % color of external links
}

%%-----------------------------------------------------------------------------------

\begin{document}

% Generate title page.
\maketitle

%%%%%%%%%%%%%%%%%%%%%%%%%%%%%%%%%%%%%%%%%%%%%%%%%%%%%%%%%
%% STEP 3:
%%%%%%%%%%%%%%%%%%%%%%%%%%%%%%%%%%%%%%%%%%%%%%%%%%%%%%%%%
%% Write your abstract in the separate file, to be positioned here.
%% You can make several abstract pages (if you want it in different languages),
%% in which case you should also define the language of the abstract,
%% as below.

% !TEX root = ./HY-CS-main.tex
\begin{abstract}

Nykypäivänä opetuksen tukena oppilaitoksissa hyödynnetään useita erilaisia verkkopalveluita. Yhteistä jokaiselle näistä palveluista on, että ne tallentavat tietoa käyttäjän toiminnasta sekä käyttäjän suorituksien arvioinneista. Tämän tutkielman tavoitteena on perehtyä mahdollisuuksiin hyödyntää tätä tietoa oppijan oppimistuloksien sekä opetuksen parantamiseksi.

Tutkielmassa tarkastellaan oppimisen analysoimisen tarpeita ja miksi oppimisanalytiikkaa halutaan tehdä. Lisäksi tarkastellaan Moodlea datalähteenä oppimisanalytiikalle sekä tutustutaan tarkemmin kahteen erilaiseen oppimisanalytiikassa käytettyyn tietomalliin. Lopuksi tarkastellaan kuinka oppimisanalytiikan avulla voidaan tietolähteiden avulla muodostaa oppijaa sekä opetusta tukevia analyysejä.

\end{abstract}

\begin{otherlanguage}{english}
\begin{abstract}

\end{abstract}
\end{otherlanguage}


% Place ToC
%\newpage
\mytableofcontents

\mainmatter

%%%%%%%%%%%%%%%%%%%%%%%%%%%%%%%%%%%%%%%%%%%%%%%%%%%%%%%%%
%% STEP 4: Write the thesis.
%%%%%%%%%%%%%%%%%%%%%%%%%%%%%%%%%%%%%%%%%%%%%%%%%%%%%%%%%
%% Your actual text starts here. You shouldn't mess with the code above the line except
%% to change the parameters. Removing the abstract and ToC commands will mess up stuff.
%%
%% Command \include{file} includes the file of name file.tex.
%% A new page will be created at every \include command, 
%% which makes it appropriate to use it for large entities such as book chapters. Cannot be nested.
%% It is useful for a big project, as changing one of the include targets 
%% won't force the regeneration of the outputs of all the rest.
%% Alternatively, \input is a more lower level macro 
%% which simply inputs the content of the given file like it was copy&pasted there manually.

%% Change your chapter titles within these files and add the actual text in them.

\chapter{Introduction\label{intro}}


\include{Ch.20_Challenges and approaches}
\include{Ch.30_Buildingblocks for solution}
\chapter{Evaluative discussion\label{discussion}}

\chapter{Conclusions\label{conclusions}}


%%%%%%%%%%%%%%%%%%%%%%%%%%%%%%%%%%%%%%%%%%%%%%%%%%%%%%%%%
%\cleardoublepage                          %fixes the position of bibliography in bookmarks
%\phantomsection
\addcontentsline{toc}{chapter}{\bibname}  % This lines adds the bibliography to the ToC
\printbibliography

%%%%%%%%%%%%%%%%%%%%%%%%%%%%%%%%%%%%%%%%%%%%%%%%%%%%%%%%%
\backmatter
\begin{appendices}

%% A sample Appendix

\appendix{Sample Appendix\label{appendix:sample}}

You can add one or more appendices to your thesis.

%% another appendix

\appendix{Instructions for LaTex}

\section{General Setup}

In the HY-CS-main.tex file you will find the following STEPS 0--5. Below you can find related instructions.
\vspace{0.5cm}

\textbf{STEP 0 -- Access the thesis template}

\begin{itemize}
\item Import the thesis template into a new Overleaf project. The easiest way to do it is to:
\begin{itemize}
    \item Obtain a zip file of the LaTeX template from the webpage of your programme.
    \item Go to \url{https://www.overleaf.com/edu/helsinki} and login to Overleaf with your university credentials.
    \item Go to the list of your projects at \url{https://www.overleaf.com/project}, click ``New Project'' and ``Upload Project''.,  the projects under your account
    \item Then upload the zip with the template.
    \item You are now ready to write your thesis in Overleaf by editing the template, you can start by renaming the project.
\end{itemize}
\end{itemize}


{\textbf{STEP 1 -- BSc or MSc thesis?}}
\begin{enumerate}
\item Select whether your are writing BSc (tkt) or MSc (csm for CS) thesis.
\item Select your language: \texttt{finnish}, \texttt{english}, or \texttt{swedish}.
\item If you are writing MSc select your line / track.
\end{enumerate}


{\textbf{STEP 2 -- Set up your personal information}}

\begin{enumerate}
\item Specify the title of your thesis with \texttt{\textbackslash title\{\}}.
\item Specify your name to the author field with \texttt{\textbackslash author\{\}}.
\item Specify the names of your supervisors of the thesis with \texttt{\textbackslash supervisors\{\}}.
\item Specify the keywords of the thesis with \texttt{\textbackslash keywords\{\}}.
\item Specify the ACM classification terms of the thesis with \texttt{\textbackslash classification\{\}}. See \url{https://dl.acm.org/ccs} for more information.
\end{enumerate}

{\textbf{STEP 3 -- Write your abstract}}

\begin{itemize}
\item You can have the abstract in multiple languages with the \texttt{otherlanguages} environment. The example below shows how to provide an English abstract:

\begin{verbatim}
\begin{otherlanguage}{english}
\begin{abstract}
Your abstract text goes here.
\end{abstract}
\end{otherlanguage}
\end{verbatim}

\end{itemize}

{\textbf{STEP 4 -- Writing your thesis}}

\begin{enumerate}
\item There are some minimal contents and instructions below
\item Remove, or comment out, this appendix from your thesis.
\end{enumerate}

{\textbf{STEP 5 -- Set your bibliography style}}

\begin{itemize}
\item The default is Author-Year style (Einstein, 1905), but it can be easily changed to numbered [1] or alphabetical [Ein05] , as the examples of these are in comments.
\item Discuss the style to use with your supervisor.
\end{itemize}

\section{Bibliography in Latex}

The bibliography is defined in a separate \texttt{.bib} file. For this template, it is named \texttt{bibliography.bib} and includes the content show in Figure~\ref{bibexamples}.

Chapter Bibliography lists all the works that you refer to in your text. You refer to the works in the bibliography using an appropriate \emph{citation key}.
%
%This thesis template contains an example of a bibliography.


References are done using \texttt{\textbackslash citep\{einstein\}}, which generates in text a citation formatted according to the selected style \citep{einstein}, or \texttt{\textbackslash citep\{latexcompanion,knuth99\}}, which generates \citep{latexcompanion,knuth99}.
As examples of a different kinds of citations (see how these look in the Latex source), we can write \citep{einstein} to refer to the work written by \citeauthor{einstein} in \citeyear{einstein}, because the work by \citet{einstein} appears in the bilbliography included in this template.

Note that there are different possible styles for the bibliography and citation keys.
%
Consult your supervisors on the chosen style -- and once you arrive at a preferred style, use it consistently throughout the thesis.

\begin{figure}[ht]
    \centering
    \begin{scriptsize}
\begin{verbatim}
@article{einstein,
    author =       "Albert Einstein",
    title =        "{Zur Elektrodynamik bewegter K{\"o}rper}. ({German})
        [{On} the electrodynamics of moving bodies]",
    journal =      "Annalen der Physik",
    volume =       "322",
    number =       "10",
    pages =        "891--921",
    year =         "1905",
    DOI =          "http://dx.doi.org/10.1002/andp.19053221004"
}

@book{latexcompanion,
    author    = "Michel Goossens and Frank Mittelbach and Alexander Samarin",
    title     = "The \LaTeX\ Companion",
    year      = "1993",
    publisher = "Addison-Wesley",
    address   = "Reading, Massachusetts"
}

@book{knuth99,
    author    = "Donald E. Knuth",
    title     = "Digital Typography",
    year      = "1999",
    publisher = "The Center for the Study of Language and Information",
    series    = "CLSI Lecture Notes (78)"
}\end{verbatim}
\end{scriptsize}
    \caption{Examples of bibliographic reference in .bib file.}
    \label{bibexamples}
\end{figure}

%In the last reference url field the code \verb+%7E+ will translate into \verb+~+ once clicked in the final pdf.

\section{Some instructions about writing in Latex}

The following gives some superficial instructions for using this template for a Master's thesis. For guidelines on thesis writing you can consult various sources, such as university courses on scientific writing or your supervisors.

For more detailed instructions, just google, e.g., "Overleaf table positioning", and your chances of finding good info are pretty good.


\section{Figures}
Besides text, here are simple examples how you can add figures and tables in your thesis.
Remember always to refer to each figure in the main text and provide them with a descriptive caption.

Figure~\ref{fig:logo} is an example of a figure in the document (see the source about how to add them).
%Using figures is particularly useful to display plots of experimental results.

\begin{figure}[ht]
\begin{center}
\includegraphics[width=0.3\textwidth]{template/figures/HY-logo-ml.png}
\caption{University of Helsinki flame-logo for Faculty of Science.\label{fig:logo}}
\end{center}
\end{figure}

\section{Tables}

Table~\ref{table:results} gives an example of a table.
Remember always to cite the table in the main text, table captions go on top of the table.

\begin{table}[h] % h positions the table here, t! would force on top of the page, or example.
\begin{center}
\caption{Experimental results.\label{table:results}} % caption is here to make it on top
\begin{tabular}{l||l c r}
Experiment & 1 & 2 & 3 \\
\hline \hline
$A$ & 2.5 & 4.7 & -11 \\
$B$ & 8.0 & -3.7 & 12.6 \\
$A+B$ & 10.5 & 1.0 & 1.6 \\
\hline
%
\end{tabular}
\end{center}
\end{table}



%% yet another appendix

\appendix{Tutkielmapohjan käyttöohjeet}
\label{appendix:instructions_finnish}

\section{Ensiaslkeleet}

\texttt{HY-CS-main.tex} tiedosto sisältää viisi askelta STEPS 0--5. Alla on kuvattu, mitä nämä askeleet tarkoittavat ja miten niitä seuraamalla luot pohjan tutkielmallesi.
\vspace{0.5cm}

\textbf{STEP 0 -- Kopioi tutkielmapohja}

\begin{itemize}
\item Hae tutkielmapohja uuteen Overleaf-projektiin. Tämä käy helpoiten seuraavasti:
\begin{itemize}
    \item Lataa Latex-pohjan zip-tiedosto koulutusohjelman sivuilta.
    \item Mene osoitteeseen \url{www.overleaf.com/edu/helsinki} ja kirjaudu Overleafiin yli\-opiston tunnuksillasi.
    \item Overleafissa (\url{https://www.overleaf.com/project}), klikkaa ``New Project'' and ``Upload Project''.
    \item Valitse lataamasi tutkielmapohjan zip-tiedosto.
    \item Nyt voit lähteä kirjoittamaan tutkielmaasi suoraan pohjaan, voit aloittaa esim. vaihtamalla projektin nimen.
\end{itemize}
\end{itemize}


{\textbf{STEP 1 -- BSc vai MSc tutkielma?}}
\begin{enumerate}
\item Valitse (tiedostossa \texttt{HY-CS-main.tex}) oletko tekemässä BSc (tkt) vai MSc (csm tietojenkäsittely) tutkielmaa.
\item Valitse kieli jolla kirjoitat tutkielman: \texttt{finnish}, \texttt{english} tai \texttt{swedish}.
\item Jos olet kirjoittamassa maisterintutkielmaa, valitse linja/opintosuunta.
\end{enumerate}


{\textbf{STEP 2 -- Aseta henkilökohtaiset tietosi}}

\begin{enumerate}
\item Kirjoita alustava otsikko tutkielmallesi: \texttt{\textbackslash title\{\}}.
\item Kirjoita oma nimesi kohtaan \texttt{\textbackslash author\{\}}.
\item Lisää ohjaajien nimet \texttt{\textbackslash supervisors\{\}}.
\item Määrittele avainsanat \texttt{\textbackslash keywords\{\}}.
\item Määritä tutkielmasi ACM luokittelutermit \texttt{\textbackslash classification\{\}}. Ks. lisätietoa: \url{https://dl.acm.org/ccs}.
\end{enumerate}

{\textbf{STEP 3 -- Kirjoita tiivistelmä}}

%\begin{itemize}
%\item
Voit kirjoittaa tiivistelmän (koko tiivistelmäsivu) eri kielillä \texttt{otherlanguages}-ym\-pä\-ris\-tön avulla. Alla esimerkki jolla kirjoitat englanninkielisen tiivistelmän muulla kuin englannin kielellä kirjoitettuun tutkielmaan:

\begin{verbatim}
\begin{otherlanguage}{english}
\begin{abstract}
Your abstract text goes here.
\end{abstract}
\end{otherlanguage}
\end{verbatim}

%\end{itemize}

{\textbf{STEP 4 -- Kirjoita tutkielma}}

\begin{enumerate}
\item Kirjoittamisesta Latexilla löydät hieman ohjeita alempaa.
\item Poista tämä liite ja muu ohjeistus tutkielmastasi, esim. kommentoimalla.
\end{enumerate}

{\textbf{STEP 5 -- Aseta kirjallisuuslähdeluettelon tyyli}}

\begin{itemize}
\item Oletustyylin tekijä-vuosi, eli (Einstein, 1905), voit vaihtaa viittaustyylin (tiedostossa \texttt{HY-CS-main.tex}) helposti (eri mallit kommentoituna) esim. numeroituun [1], tai aakkostyyliin [Ein05].
Lisää ohjeita liittyen viittaustyylin säätämiseen {Bib}\TeX issä löytyy verkosta: \url{https://ctan.org/pkg/biblatex}
\item Sovi käytettävä tyyli ohjaajasi kanssa.
\end{itemize}

\section{Kirjallisuusviitteet Latexissa}

Kirjallisuuslähteet ylläpidetään erillisessä .bib-tiedostossa. Tässä tutkielmapohjassa käy\-te\-tyt kirjallisuuslähteet, joista esimerkkejä kuvassa~\ref{bibexamples-fi}, löytyvät tiedostosta\newline \texttt{bibliography.bib}.

\begin{figure}[ht]
    \centering
    \begin{scriptsize}
\begin{verbatim}

@article{einstein,
    author =       "Albert Einstein",
    title =        "{Zur Elektrodynamik bewegter K{\"o}rper}. ({German})
        [{On} the electrodynamics of moving bodies]",
    journal =      "Annalen der Physik",
    volume =       "322",
    number =       "10",
    pages =        "891--921",
    year =         "1905",
    DOI =          "http://dx.doi.org/10.1002/andp.19053221004"
}

@book{latexcompanion,
    author    = "Michel Goossens and Frank Mittelbach and Alexander Samarin",
    title     = "The \LaTeX\ Companion",
    year      = "1993",
    publisher = "Addison-Wesley",
    address   = "Reading, Massachusetts"
}

@book{knuth99,
    author    = "Donald E. Knuth",
    title     = "Digital Typography",
    year      = "1999",
    publisher = "The Center for the Study of Language and Information",
    series    = "CLSI Lecture Notes (78)"
}
\end{verbatim}
\end{scriptsize}
    \caption{Esimerkkejä kirjallisuuslähteiden kuvaamisesta .bib-tiedostossa.}
    \label{bibexamples-fi}
\end{figure}

Viitteet kirjallisuuslähteisiin muodostetaan komennolla \texttt{\textbackslash citep\{einstein\}}, josta generoituu tekstiin valitun viittaustyylin mukaisesti muotoiltu viite \citep{einstein}, tai \texttt{\textbackslash citep\{latexcompanion,knuth99\}}, josta tekstiin puolestaan generoituu \citep{latexcompanion,knuth99}.
Voit esimerkiksi kirjoittaa \citep{einstein} viitataksesi julkaisuun, jonka on kirjoittanut \citeauthor{einstein} vuonna \citeyear{einstein}, kun vain lähde \citet{einstein} on oikein lisättynä kirjallisuuslähdetiedostossa (katso miltä nämä näyttävät Latex lähdekoodissa).

Tekstissä viitatut kirjallisuuslähteet tulevat automaattisesti viiteluetteloon. Kirjallisuuslähteiden tietojen oikeellisuus ja yhdenmukaisuus .bib-tiedostossa vaikuttavat luonnollisesti siihen, miten tiedot tutkielmassa näyttäytyvät. Tämä on syytä huomioida, sillä esim.\ verkosta valmiiksi {Bib\TeX} muodossa löytyvien tietojen täydellisyyten tai samanmuotoisuuteen ei pidä sokeasti luottaa.


Keskustele viittaustyylin valinnasta ohjaajan kanssa.
%Joitain vaihtoehtoja on osoitteessa:\\
%\url{https://www.overleaf.com/learn/latex/Biblatex_bibliography_styles}.
%\url{https://www.sharelatex.com/learn/Bibtex_bibliography_styles}.

\section{Joitain ohjeita Latexilla kirjoittamiseen}

Seuraavassa on joitain ohjeita tämän tutkielmapohjan käyttöön maisterintutkielmassa. Kirjoittamisohjeita löytyy useasta eri lähteestä. Voit esimerkiksi tutustua kandidaatintutkielman ohjeisiin.
Ohjaajan kanssa on hyvä keskustella aikaisessa vaiheessa työn rakenteesta.

Yksityiskohtaisia ohjeita Latexin käyttämäsestä saa parhaiten hakemalla verkosta, esim. haku englanniksi "Overleaf table positioning" tuottaa oletettavasti aika toimivan vastauksen.

\section{Kuvat}
Kuva~\ref{fig:logo-fi} toimii esimerkkinä kuvan lisäämisestä työhön (katso tarkemmin mallia Latex lähdekoodista). Muista myös viitata jokaiseen kuvaan tekstissä.

\begin{figure}[ht] % remove [h!] for automatic placement, which is probably better for a thesis with more text on page
\centering
\includegraphics[width=0.3\textwidth]{template/figures/HY-logo-ml.png}
\caption{Helsingin yliopiston logo matemaattis-luonnontieteellisen tiedekunnan värein.\label{fig:logo-fi}}
\end{figure}

\newpage % just to keep the table on the same page with the short piece of text
\section{Taulukot}

Taulukossa~\ref{table:results-fi} on esimerkki kokeellisten tulosten raportoinnista taulukkona. Muista myös viitata jokaiseen taulukkoon tekstissä.
\begin{table}[ht]
\centering
\caption{Kokeelliset tulokset.\label{table:results-fi}}
\begin{tabular}{l||l c r}
Koe & 1 & 2 & 3 \\
\hline \hline
$A$ & 2.5 & 4.7 & -11 \\
$B$ & 8.0 & -3.7 & 12.6 \\
$A+B$ & 10.5 & 1.0 & 1.6 \\
\hline
%
\end{tabular}
\end{table}






\end{appendices}
%%%%%%%%%%%%%%%%%%%%%%%%%%%%%%%%%%%%%%%%%%%%%%%%%%%%%%%%%

\end{document}
